\documentclass[10pt, letterpaper]{article}

\usepackage[
    ignoreheadfoot,
    top=1.5 cm,
    bottom=1.5 cm,
    left=1.5 cm,
    right=1.5 cm,
    footskip=1.0 cm,
]{geometry}
\usepackage{titlesec}
\usepackage{tabularx}
\usepackage{array}
\usepackage[dvipsnames]{xcolor}
\definecolor{primaryColor}{RGB}{0, 0, 0}
\usepackage{enumitem}
\usepackage{fontawesome5}
\usepackage{amsmath}
\usepackage[
    pdftitle={Mohamed Bnshi - Software Engineer},
    pdfauthor={Mohamed Bnshi},
    colorlinks=true,
    urlcolor=primaryColor
]{hyperref}
\usepackage[pscoord]{eso-pic}
\usepackage{calc}
\usepackage{bookmark}
\usepackage{lastpage}
\usepackage{changepage}
\usepackage{paracol}
\usepackage{ifthen}
\usepackage{needspace}
\usepackage{iftex}

\ifPDFTeX
    \input{glyphtounicode}
    \pdfgentounicode=1
    \usepackage[T1]{fontenc}
    \usepackage[utf8]{inputenc}
    \usepackage{lmodern}
\fi

\usepackage{charter}

\raggedright
\AtBeginEnvironment{adjustwidth}{\partopsep0pt}
\pagestyle{empty}
\setcounter{secnumdepth}{0}
\setlength{\parindent}{0pt}
\setlength{\topskip}{0pt}
\setlength{\columnsep}{0.15cm}
\pagenumbering{gobble}

\titleformat{\section}{\needspace{4\baselineskip}\bfseries\large}{}{0pt}{}[\vspace{1pt}\titlerule]
\titlespacing{\section}{-1pt}{0.3cm}{0.2cm}

\renewcommand\labelitemi{$\vcenter{\hbox{\small$\bullet$}}$}
\newenvironment{highlights}{
    \begin{itemize}[topsep=0.10cm, parsep=0.10cm, partopsep=0pt, itemsep=0pt, leftmargin=10pt]
}{
    \end{itemize}
}

\newenvironment{onecolentry}{
    \begin{adjustwidth}{0cm}{0cm}
}{
    \end{adjustwidth}
}

\let\hrefWithoutArrow\href

\begin{document}

    % ===== HEADER =====
    \begin{center}
        {\fontsize{25pt}{25pt}\selectfont Mohamed Bnshi}

        \vspace{5pt}

        \normalsize
        Karlskrona \quad
        \hrefWithoutArrow{mailto:mohamedbnshi@gmail.com}{mohamedbnshi@gmail.com} \quad
        079-018 44 41 \quad
        \hrefWithoutArrow{https://www.linkedin.com/in/mohamedbnshi/}{LinkedIn} \quad
        \hrefWithoutArrow{https://github.com/Mobn23/}{GitHub} \quad
        \hrefWithoutArrow{https://portfolio-delta-virid-21.vercel.app/}{Portfolio}
    \end{center}

    % ===== PROFIL =====
    \section{Profil}
    \begin{onecolentry}
        Programvaruutvecklare med erfarenhet av fullstack-utveckling, DevOps och AI-drivna system. Arbetar med CI/CD, containerisering och LLM-baserade agenter i industriell miljö genom samarbete med Ericsson. Stark grund inom säker mjukvaruutveckling och DevSecOps. Söker roll inom Software Engineering, fullstack-utveckling, DevOps eller IT-konsulting.
    \end{onecolentry}

    % ===== KOMPETENSÖVERSIKT =====
    \section{Kompetensöversikt}
    \begin{onecolentry}
        \begin{highlights}
            \item Fullstack-utveckling (React, Node.js, Python, Flask, PHP)
            \item DevOps \& CI/CD (Docker, Kubernetes, GitHub Actions, Azure DevOps, Ansible)
            \item AI-drivna system och automatisering (LLM, AI-agenter, MCP)
            \item Säker mjukvaruutveckling \& DevSecOps (OWASP, sårbarhetsanalys, threat modeling)
        \end{highlights}
    \end{onecolentry}

    % ===== UTBILDNING =====
    \section{Utbildning}
    \begin{onecolentry}
        \textbf{Blekinge Tekniska Högskola (BTH) | Karlskrona} \\
        Kandidatexamen i Programvaruteknik (Webbprogrammering, 180 hp) | Pågående (Termin 6)
        \begin{highlights}
            \item Huvudområde: Programvaruteknik
            \item Fokus: Fullstack-utveckling, Systemarkitektur och Databasteknik
        \end{highlights}
        Examen: Augusti 2026 | Tillgänglig: Sommar 2026
    \end{onecolentry}

    % ===== TEKNISKA FÄRDIGHETER =====
    \section{Tekniska Färdigheter}
    \begin{onecolentry}
        \textbf{Programspråk:} JavaScript, TypeScript, Python, PHP, SQL, Bash \\[0.15cm]
        \textbf{Ramverk \& Bibliotek:} React, Node.js, Express.js, Flask, Symfony 7, Astro, Tailwind CSS, Web Components \\[0.15cm]
        \textbf{Databaser:} MySQL, PostgreSQL, SQLite, MongoDB \\[0.15cm]
        \textbf{APIer:} RESTful APIs, GraphQL \\[0.15cm]
        \textbf{ORM/ODM:} Doctrine (PHP), Mongoose (MongoDB) \\[0.15cm]
        \textbf{DevOps \& Verktyg:} Docker, Kubernetes, Azure DevOps, GitHub Actions, Jenkins, Gerrit, Ansible, Git, Linux/Unix \\[0.15cm]
        \textbf{Säkerhet:} OWASP, sårbarhetsanalys (CVE, CVSS), SBOM, threat modeling, container-scanning \\[0.15cm]
        \textbf{Regelverk:} GDPR-medveten utveckling och hantering av känslig data
    \end{onecolentry}

    % ===== AKADEMISKT SAMARBETE =====
    \section{Akademiskt Samarbete}

    \begin{onecolentry}
        \textbf{Ericsson-BTH | Software Engineer (Student Position) -- AI-driven Systems} \hfill \textit{Pågående}
        \\[0.15cm]
        \textbf{Team av 11 ingenjörer | Agilt arbetssätt}
        \begin{highlights}
            \item Designar och implementerar MCP-servrar för att tillhandahålla kontextuell data till LLM-baserade agenter
            \item Utvecklar MCP-klienter som möjliggör återanvändbara och skalbara AI-drivna arbetsflöden
            \item Designar automatiserade feedback-loopar som analyserar CI-fel och genererar korrigerande Gerrit-patchar, vilket minskar manuellt felsökningsarbete
            \item Containeriserar agenttjänster och testmiljöer med Docker och Kubernetes för tillförlitlighet och reproducerbarhet
            \item Bidrar till att bygga robusta och återanvändbara AI-drivna automationssystem i enterprise-miljöer med fokus på skalbarhet, tillförlitlighet och minskat manuellt arbete
        \end{highlights}
        \textbf{Tekniker:} MCP, LLM-integration, AI-agenter, Event-Driven Architecture, Docker, Kubernetes, CI/CD, Gerrit, Jenkins
    \end{onecolentry}

    \vspace{0.3cm}

    \begin{onecolentry}
        \textbf{Examensarbete | BTH i samarbete med Ericsson} \hfill \textit{Pågående} \\
        Operationalisering av kontextmedveten sårbarhetsanalys i CI/CD-pipelines inom DevSecOps
        \begin{highlights}
            \item Utvecklar proof-of-concept för kontextmedveten sårbarhetsprioritering i automatiserade CI/CD-pipelines
            \item Aggregerar och normaliserar resultat från sårbarhetsskannrar till enhetligt format med deduplicering
            \item Analyserar och prioriterar sårbarheter baserat på systemkontext (exponering, säkerhetskontroller, hotmodellering)
            \item Genomförs i samarbete med Ericsson i industriell miljö
            \item Undersöker hur threat modeling kan bidra med systemkontext för förbättrad sårbarhetsprioritering i en industriell DevSecOps-miljö
            \item Designar och bygger automatiserade pipeline-flöden för beslutsfattande i CI/CD-processer
        \end{highlights}
        \textbf{Teknologier:} CI/CD, DevSecOps, container-scanning, sårbarhetsanalys, pipeline-automation, threat modeling
    \end{onecolentry}

    \vspace{0.3cm}

    \begin{onecolentry}
        \textbf{Trafikverket -- Studentutvecklare, DevOps} \hfill \textit{Jan 2025 -- Jun 2025} \\

            Akademiskt samarbete (sekretessbelagt) | Team av 7 studenter \\
        \begin{highlights}
            \item DevOps Engineer med ansvar från kravanalys och kundkontakt till implementation och deployment
            \item Implementerade automatiserad ``Vulnerability Checker'' i teamets CI/CD-pipeline för stärkt säkerhet i leveranskedjan
            \item Genomförde sårbarhetsanalyser med fokus på CVE, CVSS och SBOM-hantering
            \item Huvudpresentatör vid slutpresentation mot kund - demonstrerade teknisk lösning och affärsnytta
        \end{highlights}
    \end{onecolentry}

    % ===== PROJEKT =====
    \section{Projekt (Urval)}

    \begin{onecolentry}
        \textbf{Microservices Uthyrningssystem för Elsparkcyklar} \hfill \textit{BTH, 2025}
        \begin{highlights}
            \item Fullstack-utveckling av realtidssystem baserat på microservices-arkitektur
            \item Docker för containerhantering, CI-testning via Scrutinizer
            \item Bidrog i planeringen av hela teknikstacken
        \end{highlights}
    \end{onecolentry}

    \vspace{0.2cm}

    \begin{onecolentry}
        \textbf{MoveOut -- Flyttsystem} \hfill \textit{BTH, 2024}
        \begin{highlights}
            \item Monolitisk webbapplikation utvecklad enligt Scrum-metodik med MySQL och ren SQL
            \item Individuellt ansvarig för planering och fullstack-implementering
        \end{highlights}
    \end{onecolentry}

    \vspace{0.2cm}

    \begin{onecolentry}
        \textbf{DevOps-projekt -- CI/CD och molndriftsättning} \hfill \textit{BTH, 2025--2026}
        \begin{highlights}
            \item Designade och implementerade CI/CD-pipeline för bygg, test och driftsättning av applikation
            \item Arbetade med containerisering (Docker) och molnbaserad miljö (Azure)
            \item Automatiserade konfiguration och deployment med Infrastructure as Code-principer
            \item Använde versionshantering och DevOps-arbetsflöden för kontinuerlig leverans
            \item Erfarenhet av testautomation, staging-miljöer och övervakning
        \end{highlights}
        \textbf{Tekniker:} Docker, Kubernetes, CI/CD, Azure, Ansible, Git, Linux
    \end{onecolentry}

    \vspace{0.2cm}

    \begin{onecolentry}
        \textbf{Övriga webbprojekt:} Symfony 7 MVC-app med Doctrine ORM \textbullet{} SPA med Web Components \& GraphQL \textbullet{} Chatbot i Python
    \end{onecolentry}

    % ===== SPRÅK =====
    \section{Språk}
    \begin{onecolentry}
        Svenska -- Flytande \quad | \quad Engelska -- Flytande \quad | \quad Arabiska -- Modersmål
    \end{onecolentry}

\end{document}
